% !TeX program = xelatex
% !TeX encoding = UTF-8 Unicode
\documentclass[UTF8,openany,12pt]{ctexbook}

\usepackage[a4paper]{geometry}
\usepackage{graphicx}
\usepackage{booktabs}
\usepackage{longtable}
\usepackage{array}
\usepackage{tabularx}
\usepackage{amsmath,amssymb}
\usepackage{float}
\usepackage{caption}
\usepackage{enumitem}
\usepackage{tocloft}
\usepackage{hyperref}
\usepackage{setspace}
\usepackage{xcolor}
\usepackage[numbers,sort&compress]{natbib}

\geometry{
    top=2.54cm,
    bottom=2.54cm,
    left=3.17cm,
    right=3.17cm,
    headheight=1.2cm,
    headsep=0.4cm,
    footskip=1.6cm
}

\graphicspath{{../img/}}
\hypersetup{hidelinks}
\setlength{\parindent}{2em}
\setlength{\parskip}{0pt}
\linespread{1.35}
\raggedbottom

\DeclareCaptionFont{fsuwuhao}{\zihao{5}\songti}
\captionsetup{font=fsuwuhao,labelfont=normalfont,labelsep=space}
\captionsetup[figure]{justification=centering}
\captionsetup[table]{justification=centering,singlelinecheck=false,skip=6pt}

\setlength{\tabcolsep}{4.5pt}
\renewcommand{\arraystretch}{1.22}
\setcounter{secnumdepth}{3}
\setcounter{tocdepth}{2}
\pagestyle{plain}

\ctexset{
    chapter = {
        number = \arabic{chapter},
        format = \centering\heiti\zihao{-2},
        beforeskip = 24pt,
        afterskip = 18pt
    },
    section = {
        format = \raggedright\heiti\zihao{3},
        beforeskip = 18pt,
        afterskip = 6pt
    },
    subsection = {
        format = \raggedright\heiti\zihao{4},
        beforeskip = 12pt,
        afterskip = 6pt
    }
}

\newcommand{\keywords}[1]{\par\noindent{\heiti 关键词：}#1}
\newcommand{\engkeywords}[1]{\par\noindent{\bfseries Key words:}#1}
\newcommand{\formalmetric}[1]{\texttt{#1}}

\title{面向高召回门控的地下排水管道缺陷 Stage-1 正式容量扫描与选模研究}
\author{}
\date{2026年4月}

\begin{document}
\zihao{-4}\songti

\frontmatter

\chapter*{摘要}
\addcontentsline{toc}{chapter}{摘要}
地下排水管道缺陷识别在工程部署中同时面临三类约束：正常帧远多于缺陷帧、小样本和细粒度异常容易被背景纹理淹没，以及默认阈值下的分类指标并不能直接反映系统在高召回筛查中的减负收益。针对这一问题，本文将当前研究范围严格限定在 stage-1，并把第一阶段重新定义为 recall-constrained 的 Normal/Abnormal gate，而不是默认阈值二分类任务。本文的核心目标不是提出新的 detector，而是在已完成的 formal capacity scan 证据基础上，建立一套可复现、可审计的 stage-1 正式选模协议。

在统一的 formal protocol 下，本文分别完成了两套容量扫描：其一是 direct binary gate five-model capacity scan，其二是 source-side six-class capacity scan。前者使用 \formalmetric{Spec@R99.5}、\formalmetric{Spec@R99.0}、\formalmetric{Prec@R99.0} 与 \formalmetric{PTR@R99.0} 作为正式排序准则，并对每个 epoch checkpoint 执行温度校准与 operating-point 评估；后者仅作为 source-side representation 的辅助参考，按 Accuracy、AUROC 与 AUPRC 排序。实验结果表明，在 direct binary gate 的正式视图下，\texttt{yolo11m-cls} 取得最优综合表现，\texttt{yolo11x-cls} 为第二对照模型，而 \texttt{yolo11l-cls} 退居第三。与此同时，source six-class 视图得到的排序为 \texttt{x > m > l > s > n}，与 binary gate 视图的 \texttt{m > x > l > s > n} 并不一致。

进一步地，五个 binary gate 模型全部表现出 trainer \formalmetric{top1-best} 与 formal gate-best 的错位，说明 \formalmetric{top1} 虽然仍可作为训练健康信号，但并不对应高召回门控真正关心的 operating-point 目标，因而不应再被用作 stage-1 的正式 early-stop 或 checkpoint-selection 依据。当前结果由此建立了三点正式证据：direct binary gate 是第一阶段主裁判视图；source six-class 只承担辅助参考角色；后续 stage-1 增强实验应以 \texttt{yolo11m-cls} 为主模型，并以 \texttt{yolo11x-cls} 作为容量对照模型。基于这些结果，本文为后续的 hard-negative 回流、训练侧样本组织优化以及后续系统级研究奠定了统一的 stage-1 证据基础。

\keywords{地下排水管道；高召回门控；容量扫描；温度校准；检查点选模；CCTV 图像}

\chapter*{Abstract}
\addcontentsline{toc}{chapter}{Abstract}
This manuscript focuses exclusively on stage-1 model selection for sewer-defect screening under a high-recall operating regime. Rather than treating stage-1 as a default-threshold binary classifier, the study formalizes it as a recall-constrained Normal/Abnormal gate whose primary objective is normal filtering under fixed recall anchors. The contribution of the present manuscript is therefore not a new detector, but a reproducible and auditable stage-1 protocol built on completed formal capacity scans.

Two formal capacity scans are reported. The primary view is a direct binary-gate scan over five YOLO11 classification backbones, ranked by \formalmetric{Spec@R99.5}, \formalmetric{Spec@R99.0}, \formalmetric{Prec@R99.0}, and \formalmetric{PTR@R99.0}, with temperature scaling and operating-point evaluation applied to every epoch checkpoint. The auxiliary view is a six-class source-side scan ranked by Accuracy, AUROC, and AUPRC. The results show that \texttt{yolo11m-cls} is the formal leader under the binary gate objective, followed by \texttt{yolo11x-cls}, whereas the auxiliary six-class view ranks the models as \texttt{x > m > l > s > n}. This mismatch indicates that source-side representation quality cannot replace gate-aware checkpoint selection for stage-1.

The experiments further show that trainer-selected \formalmetric{top1-best} checkpoints do not align with gate-selected checkpoints for any of the five binary-gate models. This evidence supports a key methodological claim of the paper: trainer-side \formalmetric{top1} remains useful for monitoring optimization health, but it should not be used as the official early-stopping or model-selection criterion for a high-recall gate. The completed capacity scans therefore establish a formal stage-1 mainline centered on \texttt{yolo11m-cls}, with \texttt{yolo11x-cls} as the second control model, and provide the evidence base for subsequent hard-negative and training-side enhancement experiments.

\engkeywords{sewer inspection, high-recall gate, capacity scan, temperature scaling, checkpoint selection, CCTV imagery}

\tableofcontents

\mainmatter

\chapter{绪论}

\section{研究背景与问题提出}
地下排水管网是城市基础设施体系中的关键组成部分，其运行状态直接关系到城市防涝、道路安全以及水环境治理 \cite{moradi2019review,haurum2020survey,ma2026stateoftheart}。当前工程实践仍主要依赖 CCTV 巡检和人工复核，但海量正常帧、复杂背景纹理以及跨设备成像差异，使人工判读在效率、一致性和误报控制方面面临明显瓶颈 \cite{haurum2020survey,cheng2018deep}。因此，如何在保持缺陷高召回的前提下尽可能减少进入后续人工复核或后级检测的正常样本，已经成为 sewer inspection 智能化中的核心问题之一。

传统研究通常把第一阶段视作普通分类任务，默认以 Accuracy 或 \formalmetric{top1} 作为早停和选模依据。然而，在实际两阶段工作流中，stage-1 的职责并不是在默认阈值下取得最高整体分类分数，而是在固定高召回约束下最大化正常样本拦截。因此，默认阈值分类语义与高召回门控语义之间存在明显错位。本文据此将 stage-1 明确重定义为 recall-constrained gate，并围绕这一目标重新组织容量扫描、校准评估和 checkpoint 选择规则。

\section{相关研究现状}
\subsection{排水管道缺陷视觉识别}
排水管道缺陷识别经历了由人工特征到深度视觉模型的技术演进。早期工作以 HOG、SVM 和规则驱动候选区域为主，已经证明 CCTV 图像中的管道异常具有一定可计算性，但方法对场景变化较为敏感 \cite{halfawy2014hogsvm}。随着深度学习的发展，研究开始转向图像级分类与目标检测。Kumar 等利用深度卷积网络完成 sewer CCTV 缺陷分类，Cheng 等进一步将深度模型用于管道缺陷检测，表明数据驱动表征能够显著提升识别性能 \cite{kumar2018cnn,cheng2018deep}。

\subsection{YOLO 路线与 sewer 任务适配}
YOLO 系列因具备较好的实时性和工程成熟度，已成为管道缺陷检测中的常见基础路线 \cite{redmon2016yolo}。在 sewer 场景中，研究者围绕特征增强、轻量化部署与场景适配提出了多种改进策略，例如改进 YOLOv5 的管道缺陷检测模型以及轻量化 YOLOv5-Sewer 等 \cite{wang2023improvedyolov5,zhao2024yolov5sewer}。这些工作证明，YOLO 路线具有良好的工程基础，但也表明不同任务目标之间不能简单共用同一评价口径：用于目标检测的 detector 精度并不直接回答高召回门控下的 normal filtering 问题。

\subsection{公开数据、校准与迁移视角}
Sewer-ML 数据集的发布为 sewer defect classification 提供了更统一的 benchmark 条件，也使多尺度容量比较与 source-side representation 分析成为可能 \cite{haurum2021sewerml}。另一方面，温度校准已被证明是改善神经网络输出分数可信度的低成本方式，在不改变样本排序的前提下，有助于建立更稳定的 operating-point 比较基础 \cite{guo2017calibration}。此外，迁移学习在 sewer defect detection 中也已显示出缓解目标域标注不足的价值 \cite{situ2023transfer}。然而，就当前已完成证据而言，最需要优先解决的问题并不是 detector 结构本身，而是 stage-1 的正式任务边界与正式选模规则。

\section{研究目标与论文贡献}
本文当前版本只回答一个问题：在统一的 stage-1 formal protocol 下，哪一种 YOLO11-cls 容量更适合作为高召回 gate 的正式起点。围绕这一目标，本文形成以下三点贡献：
\begin{enumerate}[label=\arabic*.]
    \item 将 stage-1 重新界定为 recall-constrained high-recall gate，并给出一套可复现的 formal checkpoint-selection 协议，即对每个 epoch checkpoint 执行温度校准、阈值扫描和 operating-point 排序，而不再依赖 trainer-side \formalmetric{top1}。
    \item 完成两套 formal capacity scan，并用 direct binary gate 与 source six-class 的排序差异，说明 source-side representation 视图不能替代 stage-1 gate-aware 选模视图。
    \item 基于完整的 checkpoint 级证据表明，五个 binary gate 模型全部存在 \formalmetric{top1-best} 与 gate-best 的系统性错位，从而支持“\formalmetric{top1} 不应作为 stage-1 正式选模标准”这一结论。
\end{enumerate}

\section{论文组织}
全文围绕已完成的 stage-1 容量扫描结果展开。第 2 章定义 formal 指标与 checkpoint-selection 协议；第 3 章介绍数据集、容量扫描设置与评估流程；第 4 章报告 direct binary gate formal capacity scan 结果；第 5 章报告 six-class source formal capacity scan 及 cross-view rank mismatch；第 6 章讨论这些结果对 stage-1 主线选择的含义及当前边界；第 7 章总结全文并说明后续工作。

\chapter{理论基础与正式协议}

\section{stage-1 高召回门控的目标定义}
设 stage-1 的输入样本总数为 $N$，其中 abnormal 样本需要尽可能完整地保留给后续人工复核或后级处理，而 normal 样本应在尽可能不牺牲召回的前提下被前级拦截。由此，stage-1 的核心目标不是默认阈值下的总体分类正确率，而是高召回约束下的 normal filtering ability。

本文采用两个固定工作点作为正式锚点：\formalmetric{Recall$\geq$99.5\%} 与 \formalmetric{Recall$\geq$99.0\%}。在这两个约束下，正式排序规则定义为：
\begin{equation}
\formalmetric{Spec@R99.5} \succ \formalmetric{Spec@R99.0} \succ \formalmetric{Prec@R99.0} \succ \formalmetric{PTR@R99.0},
\end{equation}
其中符号“$\succ$”表示先后优先级，而非加权求和。

\section{正式指标定义}
在二分类 gate 语境中，\formalmetric{Spec@R99.5} 与 \formalmetric{Spec@R99.0} 分别表示在满足相应召回约束时的 specificity；\formalmetric{Prec@R99.0} 表示在 \formalmetric{Recall$\geq$99.0\%} 约束下的 precision。本文特别采用 \formalmetric{PTR@R99.0} 作为第四排序准则，其含义为在 \formalmetric{Recall$\geq$99.0\%} 对应阈值下的 pass-through rate，即仍被送往下一阶段的样本占全部样本的比例：
\begin{equation}
\formalmetric{PTR@R99.0} = \frac{TP + FP}{N}.
\end{equation}
其中 $TP$ 与 $FP$ 分别为进入下一阶段的 abnormal 与 normal 样本数。由于 stage-1 的目标是尽可能在不损害高召回的前提下减轻后续负担，因此 \formalmetric{PTR@R99.0} 越低越好。

\section{温度校准与 gate-aware checkpoint selection}
为避免不同 checkpoint 的分数尺度不可比，本文对每个 epoch checkpoint 采用单标量 temperature scaling \cite{guo2017calibration}。具体流程如下：
\begin{enumerate}[label=\arabic*.]
    \item 在固定的 val-cal 子集上拟合温度参数 $T$；
    \item 将该温度应用于 val-op 子集，得到 calibrated score；
    \item 在 val-op 上执行 threshold sweep，计算各候选阈值下的 recall、specificity、precision 与 pass-through rate；
    \item 按照正式排序规则选出 gate-best checkpoint，而不是直接使用 trainer 自动保存的 \texttt{best.pt}。
\end{enumerate}

这一协议意味着 calibration 已经内嵌为 formal evaluation 的组成部分，而不是独立于容量扫描之外的附加训练阶段。trainer-side 的 \formalmetric{acc/top1/loss} 仍然保留，但它们只承担训练健康监控功能，不承担正文正式选模功能。

\section{six-class source 视图的辅助地位}
除 direct binary gate 外，本文还组织了一条 six-class source 容量扫描，用于观察不同容量模型在 source-side representation 任务上的相对位置。该视图的正式排序规则为 Accuracy、AUROC 与 AUPRC。然而，这一视图只承担辅助参考角色，不应替代 direct binary gate 的正式 backbone 选择。后续章节将用实证结果说明，source six-class 排序与 direct binary gate 排序确实存在错位。

\chapter{数据集、实验设置与 formal capacity scan 设计}

\section{数据集组织}
本文使用两套 source-side 数据集来支撑当前 stage-1 研究。第一套是 direct binary gate 数据集 \texttt{sewerml\_gate2\_train7200}，共 7200 张图像，其中训练集 6480 张、验证集 720 张；类别分布为 abnormal 6000 张、normal 1200 张，训练集与验证集均保持 5:1 的类间比例。第二套是 six-class source 数据集 \texttt{sewerml\_cls6\_train7200}，同样包含 7200 张图像，并被均衡组织为六个类别，每类 1200 张，训练集与验证集分别为 6480 与 720 张。

对于 binary gate formal protocol，全部 720 张验证图像进一步按类别分层拆分为 val-cal 与 val-op 两部分，其中 val-cal 为 216 张，val-op 为 504 张；对应 abnormal/normal 的划分分别为 180/36 与 420/84。该拆分仅用于温度参数拟合与 operating-point 报告，不参与训练。

\section{模型族与统一训练设置}
所有容量扫描均以 YOLO11-cls 五个尺度模型为候选，即 \texttt{yolo11n-cls}、\texttt{yolo11s-cls}、\texttt{yolo11m-cls}、\texttt{yolo11l-cls} 和 \texttt{yolo11x-cls}。统一训练协议如下：
\begin{itemize}
    \item batch size 固定为 24；
    \item 训练轮数固定为 200；
    \item 每轮保存 checkpoint；
    \item 不使用 trainer-side \formalmetric{top1} 作为早停准则；
    \item 所有 formal 结果均由外部 evaluator 汇总。
\end{itemize}
该统一设置的目的，是把容量差异与 checkpoint-selection 规则显式分离，从而避免因为不同早停行为而掩盖真实的 gate-aware 排序。

\section{容量扫描输出材料}
为保证结果可追溯，本文对每个模型保留 run manifest、dataset manifest、best epoch manifest、all checkpoints index 以及 epoch-level summary，并在统一结果目录中生成主表、主图、附录表和 supplementary source-of-truth。由此，当前 manuscript 中的所有数字都可以追溯到对应的 formal summary，而不依赖训练日志中的零散截图或经验性口径。

\chapter{direct binary gate formal capacity scan 结果}

\section{主结果表}
\begin{table}[H]
    \centering
    \caption{direct binary gate formal capacity scan 主表}
    \label{tab:stage1-formal-gate-capacity}
    \resizebox{\textwidth}{!}{%
    \begin{tabular}{lcccccccc}
        \toprule
        Model & Best Epoch & Spec@R99.5 & Spec@R99.0 & Prec@R99.0 & PTR@R99.0 & $\tau_{99.5}$ & $\tau_{99.0}$ & Rank \\
        \midrule
        yolo11m-cls & 78  & 0.5952 & 0.6548 & 0.9348 & 0.8829 & 0.2800 & 0.3100 & 1 \\
        yolo11x-cls & 125 & 0.5476 & 0.5952 & 0.9244 & 0.8929 & 0.2900 & 0.3600 & 2 \\
        yolo11l-cls & 54  & 0.5238 & 0.5714 & 0.9205 & 0.8988 & 0.2800 & 0.3000 & 3 \\
        yolo11s-cls & 77  & 0.5238 & 0.5357 & 0.9143 & 0.9028 & 0.2500 & 0.2600 & 4 \\
        yolo11n-cls & 76  & 0.5119 & 0.5833 & 0.9224 & 0.8948 & 0.2500 & 0.3600 & 5 \\
        \bottomrule
    \end{tabular}
    }
\end{table}
\noindent\parbox{\textwidth}{\footnotesize 注：direct binary gate 是 stage-1 的正式选模视图；排序遵循 \formalmetric{Spec@R99.5 > Spec@R99.0 > Prec@R99.0 > PTR@R99.0}。}

\begin{figure}[H]
    \centering
    \includegraphics[width=0.92\textwidth]{../../../research/results/stage1_formal/capacity_scan/paper_main/figures/fig_stage1_gate_capacity_bar.png}
    \caption{direct binary gate formal capacity scan 的主指标比较}
    \label{fig:stage1-formal-gate-bar}
\end{figure}

表 \ref{tab:stage1-formal-gate-capacity} 表明，在统一的 recall-constrained 排序规则下，\texttt{yolo11m-cls} 在 \formalmetric{Spec@R99.5} 与 \formalmetric{Spec@R99.0} 两个主排序锚点上同时取得最优，并在 \formalmetric{Prec@R99.0} 与 \formalmetric{PTR@R99.0} 上保持最优综合表现。因此，当前 stage-1 formal binary gate 的主模型应固定为 \texttt{yolo11m-cls}，第二对照模型则应固定为 \texttt{yolo11x-cls}。这一结果与早期基于默认阈值印象形成的容量偏好不同，说明 stage-1 的 backbone 选择必须回到 high-recall gate objective 本身。

\section{Spec@R99.5 的 checkpoint dynamics}
\begin{figure}[H]
    \centering
    \includegraphics[width=0.90\textwidth]{../../../research/results/stage1_formal/capacity_scan/paper_main/figures/fig_stage1_gate_spec_r995_curve.png}
    \caption{五模型 \formalmetric{Spec@R99.5} 随 epoch 的演化曲线}
    \label{fig:stage1-formal-gate-spec995}
\end{figure}

\begin{figure}[H]
    \centering
    \includegraphics[width=0.90\textwidth]{../../../research/results/stage1_formal/capacity_scan/paper_main/figures/fig_stage1_gate_spec_r990_curve.png}
    \caption{五模型 \formalmetric{Spec@R99.0} 随 epoch 的演化曲线}
    \label{fig:stage1-formal-gate-spec990}
\end{figure}

图 \ref{fig:stage1-formal-gate-spec995} 与图 \ref{fig:stage1-formal-gate-spec990} 展示了 formal gate 指标随 checkpoint 变化的整体趋势。五个模型的 gate-best epoch 分别为 \texttt{76/77/78/54/125}（对应 \texttt{n/s/m/l/x}），说明 formal 最优点通常出现在训练中期或后中期，而不是默认意义上的“最后一轮”或“trainer 自动 best”。其中，\texttt{yolo11l-cls} 在较早 epoch 达到峰值，而 \texttt{yolo11x-cls} 需要更长优化轨迹才达到最优，但即便如此，\texttt{yolo11x-cls} 仍未在正式排序下超过 \texttt{yolo11m-cls}。因此，capacity 增大并不自动等价于 recall-constrained filtering 能力的稳定提升。

\section{trainer \formalmetric{top1-best} 与 gate-best 的错位}
\begin{table}[H]
    \centering
    \caption{trainer \formalmetric{top1-best} 与 gate-best 的对照}
    \label{tab:stage1-formal-top1-vs-gatebest}
    \resizebox{\textwidth}{!}{%
    \begin{tabular}{lccccccc}
        \toprule
        Model & Top1-Best Epoch & Top1-Best Top1 & Gate-Best Epoch & Gate-Best Spec@R99.5 & Gate-Best Spec@R99.0 & $\Delta$Epoch & Same \\
        \midrule
        yolo11n-cls & 120 & 0.9250 & 76  & 0.5119 & 0.5833 & 44  & No \\
        yolo11s-cls & 88  & 0.9208 & 77  & 0.5238 & 0.5357 & 11  & No \\
        yolo11m-cls & 197 & 0.9347 & 78  & 0.5952 & 0.6548 & 119 & No \\
        yolo11l-cls & 103 & 0.9306 & 54  & 0.5238 & 0.5714 & 49  & No \\
        yolo11x-cls & 120 & 0.9319 & 125 & 0.5476 & 0.5952 & -5  & No \\
        \bottomrule
    \end{tabular}
    }
\end{table}
\noindent\parbox{\textwidth}{\footnotesize 注：五个模型全部出现 \formalmetric{top1-best} 与 gate-best 不一致；\formalmetric{top1} 可作为训练健康信号，但不应作为 stage-1 的正式早停或检查点选择标准。}

\begin{figure}[H]
    \centering
    \includegraphics[width=0.92\textwidth]{../../../research/results/stage1_formal/capacity_scan/paper_main/figures/fig_stage1_gate_top1_vs_spec_dualaxis.png}
    \caption{以 \texttt{yolo11m-cls} 与 \texttt{yolo11l-cls} 为例的 \formalmetric{top1} 与 gate 指标错位}
    \label{fig:stage1-formal-gate-top1-mismatch}
\end{figure}

表 \ref{tab:stage1-formal-top1-vs-gatebest} 与图 \ref{fig:stage1-formal-gate-top1-mismatch} 共同揭示了一个关键结论：五个模型中没有任何一个满足 trainer \formalmetric{top1-best} 与 formal gate-best 一致。尤其是 \texttt{yolo11m-cls}，其 trainer 侧 \formalmetric{top1-best} 出现在第 197 轮，而 formal gate-best 出现在第 78 轮，两者相差 119 个 epoch；\texttt{yolo11l-cls} 也表现出类似错位。这一结果说明，trainer 侧的 \formalmetric{top1} 虽然仍能反映分类优化是否正常进行，但并不与 recall-constrained gate objective 对齐。因此，若仍以 \formalmetric{acc/top1} 作为 stage-1 的正式早停或选模准则，将系统性偏离高召回门控真正关心的工作点目标。

\section{对 stage-1 主线选择的含义}
从 operating-point 计数看，\texttt{yolo11m-cls} 在 \formalmetric{R99.5} 下达到 \texttt{TN/FN = 50/2}，在 \formalmetric{R99.0} 下达到 \texttt{TN/FN = 55/4}；\texttt{yolo11x-cls} 对应为 \texttt{46/2} 与 \texttt{50/4}；\texttt{yolo11l-cls} 对应为 \texttt{44/2} 与 \texttt{48/3}。这些结果表明，当前主模型切换并非由单个 checkpoint 的偶然波动导致，而是不同容量模型在 recall-constrained filtering setting 下的稳定差异。由此，后续所有 stage-1 增强实验都应围绕 \texttt{yolo11m-cls} 展开，并以 \texttt{yolo11x-cls} 作为容量对照模型。

\chapter{six-class source formal capacity scan 与 cross-view 分析}

\section{six-class source 主结果}
\begin{table}[H]
    \centering
    \caption{six-class source formal capacity scan 主表}
    \label{tab:stage1-formal-cls6-capacity}
    \begin{tabular}{lccccc}
        \toprule
        Model & Best Epoch & Accuracy & AUROC & AUPRC & Rank \\
        \midrule
        yolo11x-cls & 161 & 0.7000 & 0.9484 & 0.9885 & 1 \\
        yolo11m-cls & 130 & 0.6958 & 0.9387 & 0.9848 & 2 \\
        yolo11l-cls & 104 & 0.6944 & 0.9400 & 0.9862 & 3 \\
        yolo11s-cls & 167 & 0.6917 & 0.9406 & 0.9854 & 4 \\
        yolo11n-cls & 117 & 0.6833 & 0.9421 & 0.9870 & 5 \\
        \bottomrule
    \end{tabular}
\end{table}
\noindent\parbox{\textwidth}{\footnotesize 注：six-class source 视图只用于 source-side representation 参考，不承担 stage-1 的正式选模功能。}

\begin{figure}[H]
    \centering
    \includegraphics[width=0.88\textwidth]{../../../research/results/stage1_formal/capacity_scan/paper_main/figures/fig_stage1_cls6_capacity_bar.png}
    \caption{six-class source formal capacity scan 的五模型比较}
    \label{fig:stage1-formal-cls6-bar}
\end{figure}

在 source-side six-class 任务下，\texttt{yolo11x-cls} 为辅助参考 leader，\texttt{yolo11m-cls} 位居第二。这一结果说明，更大容量模型在源域六类表示任务上仍具有优势；但该优势并不能直接转译为 recall-constrained binary gate 下的正式最优。因此，six-class source 结果应被解释为 source-side representation 的辅助视图，而不是 direct binary gate 的替代裁判。

\section{cross-view rank mismatch}
\begin{table}[H]
    \centering
    \caption{six-class source 与 direct binary gate 的跨视图排序差异}
    \label{tab:stage1-formal-crossview-rank-gap}
    \resizebox{\textwidth}{!}{%
    \begin{tabular}{lcccc}
        \toprule
        Model & Rank in CLS6 & Rank in Binary Gate & Rank Gap & Interpretation \\
        \midrule
        yolo11x-cls & 1 & 2 & -1 & 辅助 cls6 视图相对高估该模型 \\
        yolo11m-cls & 2 & 1 & 1  & 辅助 cls6 视图相对低估该模型 \\
        yolo11l-cls & 3 & 3 & 0  & 两种视图排序一致 \\
        yolo11s-cls & 4 & 4 & 0  & 两种视图排序一致 \\
        yolo11n-cls & 5 & 5 & 0  & 两种视图排序一致 \\
        \bottomrule
    \end{tabular}
    }
\end{table}

\begin{figure}[H]
    \centering
    \includegraphics[width=0.88\textwidth]{../../../research/results/stage1_formal/capacity_scan/paper_main/figures/fig_stage1_crossview_rank_gap.png}
    \caption{six-class source 与 direct binary gate 的排序错位可视化}
    \label{fig:stage1-formal-crossview-gap}
\end{figure}

表 \ref{tab:stage1-formal-crossview-rank-gap} 与图 \ref{fig:stage1-formal-crossview-gap} 表明，source six-class 排序不能替代 direct binary gate 排序。最关键的差异集中在 \texttt{yolo11m-cls} 与 \texttt{yolo11x-cls}：前者在 source six-class 中位居第二，却在 binary gate 中成为正式主模型；后者在 source six-class 中位居第一，却在 binary gate 中退居第二。这一错位说明，source-side representation quality 虽具有参考意义，但第一阶段的正式 backbone 选择必须始终回到 recall-constrained gate objective 本身。

\chapter{讨论：capacity scan 已建立的结论与当前边界}

\section{已建立的正式证据}
截至当前版本，本文已经建立了三条正式证据链。第一，direct binary gate 是 stage-1 的正式选模视图，而不是默认阈值分类任务的延伸。第二，formal binary gate capacity scan 已经明确给出当前主模型与第二对照模型，即 \texttt{yolo11m-cls} 与 \texttt{yolo11x-cls}。第三，five-model \formalmetric{top1-best} 与 gate-best 的全不一致，证明 trainer-side \formalmetric{top1} 不应再作为 stage-1 的正式选模依据。这三点共同构成了当前 stage-1 主线切换的全部证据基础。

\section{为何当前不宣称两阶段整体优于单阶段}
本文当前只完成了 stage-1 formal capacity scan 与 auxiliary six-class capacity scan。尽管 stage-1 的高召回门控动机来自两阶段系统的工程需求，但当前尚未完成针对单阶段 detector 与两阶段系统的系统级对比实验。因此，本文不对“整体两阶段系统优于单阶段 detector”作出正式结论。当前 manuscript 的有效结论仅限于：若后续系统采用 stage-1 front-end，则其 front-end 主模型应在 direct binary gate 的正式视图下选择，而不能由默认阈值分类指标或 source six-class 结果替代决定。

\section{当前局限性与后续工作边界}
本文的边界同样需要明确。其一，当前稿件不包含 hard-negative sweep、HardMix、RCIS 或 detector 实验结果；这些实验仍在统一 formal protocol 下持续推进。其二，本文虽然证明了 \formalmetric{top1-best} 与 gate-best 的错位，但未引入更复杂的统计显著性分析，而是依赖完整 checkpoint trace 的一致性证据。其三，source six-class 视图当前只作为 auxiliary reference，其与 downstream detector initialization 的关系仍需后续实证闭环。基于这些边界，本文把后续工作限定为：先在 \texttt{yolo11m-cls} 上完成 HN sweep，并以 \texttt{yolo11x-cls} 做 cross-capacity 验证，再在新的主模型与比例固定后进入训练侧增强与系统级研究。

\chapter{结论与展望}

\section{结论}
基于当前已经完成的 formal capacity scan，本文可以稳定得到以下结论：
\begin{enumerate}[label=\arabic*.]
    \item stage-1 应被正式定义为 recall-constrained high-recall gate，而不是默认阈值二分类任务。相应地，trainer 侧的 \formalmetric{acc/top1} 只能作为训练健康信号，而不能再充当 formal best checkpoint 的依据。
    \item 在统一的 \texttt{batch=24}、\texttt{epochs=200}、每轮保存 checkpoint 与外部 gate-aware evaluation protocol 下，direct binary gate formal capacity scan 已正式确认 \texttt{yolo11m-cls} 为当前 high-recall leader，\texttt{yolo11x-cls} 为第二对照模型，\texttt{yolo11l-cls} 退居第三。
    \item source-side six-class formal capacity scan 得到的排序为 \texttt{x > m > l > s > n}，而 direct binary gate 的正式排序为 \texttt{m > x > l > s > n}。这一错位表明，source six-class 结果不能替代 direct binary gate 的正式选模。
    \item 五个 binary gate 模型全部表现出 \formalmetric{top1-best != gate-best}。这一现象直接支持 gate-aware checkpoint selection 的必要性，也说明默认 \formalmetric{top1} 早停并不适用于 stage-1 的正式论文口径。
\end{enumerate}

\section{展望}
后续工作的重点可概括为三方面。其一，在已固定的主模型 \texttt{yolo11m-cls} 上完成 formal HN sweep，并以 \texttt{yolo11x-cls} 做轻量 cross-capacity 验证，从而在同一 calibrated protocol 下确认新的最佳 hard-negative 比例。其二，在 backbone 与 HN 比例固定后，继续比较 HardMix、RCIS 等训练侧策略，评估样本组织方式是否能够在相同 recall 约束下带来净增益。其三，在 stage-1 证据链闭合后，再进入 detector 结构改进、迁移训练和系统级验证，从而在统一证据体系下讨论完整两阶段系统的整体价值。

\appendix

\chapter{附录：Stage-1 formal protocol 摘要}
\begin{table}[H]
    \centering
    \caption{Stage-1 capacity scan 的 formal protocol 摘要}
    \label{tab:stage1-capacity-protocol}
    \begin{tabularx}{\textwidth}{>{\raggedright\arraybackslash}p{3.0cm}X}
        \toprule
        项目 & 说明 \\
        \midrule
        primary selection view & direct binary gate \\
        auxiliary reference view & source six-class capacity scan \\
        binary gate ranking rule & \formalmetric{Spec@R99.5 > Spec@R99.0 > Prec@R99.0 > PTR@R99.0} \\
        six-class ranking rule & Accuracy $>$ AUROC $>$ AUPRC \\
        shared training protocol & \texttt{batch=24}, \texttt{epochs=200}, checkpoint saved every epoch \\
        checkpoint policy & every epoch checkpoint is externally evaluated; trainer \formalmetric{top1} does not determine the official winner \\
        calibration protocol & val-cal temperature fitting followed by val-op threshold sweep \\
        manuscript scope & only completed capacity-scan evidence is reported in the current paper \\
        \bottomrule
    \end{tabularx}
\end{table}

本文正文仅保留最关键的 formal 主表与主图。完整的 epoch-level summary、per-model dashboard、Top-3 checkpoint 表以及 source-of-truth manifests 已整理为 supplementary archive，用于后续补图、复核和继续追踪。

\backmatter
\begin{thebibliography}{99}
\bibitem{moradi2019review}
Moradi S, Zayed T, Golkhoo F. Review on Computer Aided Sewer Pipeline Defect Detection and Condition Assessment[J]. Infrastructures, 2019, 4(1): 10.

\bibitem{haurum2020survey}
Haurum J B, Moeslund T B. A Survey on Image-Based Automation of CCTV and SSET Sewer Inspections[J]. Automation in Construction, 2020, 111: 103061.

\bibitem{ma2026stateoftheart}
Ma D, Fang H, Wang N, et al. The State-of-the-Art in Pipe Defect Inspection with Computer Vision-Based Methods[J]. Measurement, 2026, 257: 118431.

\bibitem{halfawy2014hogsvm}
Halfawy M R, Hengmeechai J. Automated Defect Detection in Sewer Closed Circuit Television Images Using Histograms of Oriented Gradients and Support Vector Machine[J]. Automation in Construction, 2014, 38: 1--13.

\bibitem{kumar2018cnn}
Kumar S S, Abraham D M, Jahanshahi M R, et al. Automated Defect Classification in Sewer Closed Circuit Television Inspections Using Deep Convolutional Neural Networks[J]. Automation in Construction, 2018, 91: 273--283.

\bibitem{cheng2018deep}
Cheng J C P, Wang M. Automated Detection of Sewer Pipe Defects in CCTV Images Using Deep Learning Techniques[J]. Automation in Construction, 2018, 95: 155--171.

\bibitem{haurum2021sewerml}
Haurum J B, Moeslund T B. Sewer-ML: A Multi-Label Sewer Defect Classification Dataset and Benchmark[C]//Proceedings of the IEEE/CVF Conference on Computer Vision and Pattern Recognition. 2021: 13456--13467.

\bibitem{redmon2016yolo}
Redmon J, Divvala S, Girshick R, Farhadi A. You Only Look Once: Unified, Real-Time Object Detection[C]//Proceedings of the IEEE Conference on Computer Vision and Pattern Recognition. 2016: 779--788.

\bibitem{guo2017calibration}
Guo C, Pleiss G, Sun Y, Weinberger K Q. On Calibration of Modern Neural Networks[C]//Proceedings of the 34th International Conference on Machine Learning. 2017: 1321--1330.

\bibitem{situ2023transfer}
Situ Z, Teng S, Feng W, et al. A Transfer Learning-Based YOLO Network for Sewer Defect Detection in Comparison to Classic Object Detection Methods[J]. Developments in the Built Environment, 2023, 15: 100191.

\bibitem{wang2023improvedyolov5}
Wang T, Li Y, Zhai Y, et al. A Sewer Pipeline Defect Detection Method Based on Improved YOLOv5[J]. Processes, 2023, 11(8): 2508.

\bibitem{zhao2024yolov5sewer}
Zhao X, Xiao N, Cai Z, et al. YOLOv5-Sewer: Lightweight Sewer Defect Detection Model[J]. Applied Sciences, 2024, 14(5): 1869.
\end{thebibliography}

\end{document}
